\SourceLeaf{408}%
um $\frac13$ des sphärischen Excesses verminderten Winkel proportional sind, auf krumme Flächen \textls{jeder} Art (wo die Vertheilung \textls{ungleich} geschehen muss), welche ich der Materie nach schon seit vielen Jahren besessen, aber noch nicht zur möglichen Mittheilung an andere entwickelt hatte, jetzt in eine überaus elegante Gestalt gebracht. \mbox{. . . . . . .}
\par\Needspace{15\baselineskip}
\begin{center}[II.]\end{center}
\begin{center}[1.]\end{center}
\begin{center}EINFACHSTER BEWEIS VON LEGENDRES LEHRSATZ DIE SPHÄRISCHEN DREIECKE BETREFFEND.\end{center}
\begin{center}\begin{tabular}{ll}
$3\omega$ & Sphärischer Excess\\
$A+\omega,B+\omega,C+\omega$ & Winkel\\
$a,b,c$ & Seiten\\
\multicolumn{2}{c}{$A+B+C=180^0.$}
\end{tabular}\end{center}
\noindent\hspace*{1.2em}Man hat dann nach aller Schärfe
\[
\begin{aligned}
\sin\tfrac12 a^2&=\frac{\sin\tfrac32\omega.\sin(A-\tfrac12\omega)}{\sin(B+\omega).\sin(C+\omega)}\\
\cos\tfrac12 a^2&=\frac{\sin(B-\tfrac12\omega).\sin(C-\tfrac12\omega)}{\sin(B+\omega).\sin(C+\omega)}
\end{aligned}
\]
\[
\frac{\sin\frac12 a^6}{\cos\frac12 a^2}=\frac{\sin\frac32\omega^3.\sin(A+\omega)^2.\sin(A-\frac12\omega)^4}{\begin{gathered}\sin(A+\omega)^2.\sin(A-\frac12\omega).\sin(B+\omega)^2\\.\sin(B-\frac12\omega).\sin(C+\omega)^2.\sin(C-\frac12\omega)\end{gathered}}
\]
\[
\frac{\sin\frac12 a^3}{\cos\frac12 a}.\frac{\cos\frac12 b}{\sin\frac12 b^3}=\frac{\sin(A+\omega).\sin(A-\frac12\omega)^2}{\sin(B+\omega).\sin(B-\frac12\omega)^2}
\]
\[
\begin{aligned}
\frac ab=\frac{\sin A}{\sin B}\mkern6mu{}^3\mkern-4mu\BareRoot{\Biggl(}&\frac{a^3\cos\frac12 a}{8\sin\frac12 a^3}.\frac{8\sin\frac12 b^3}{b^3\cos\frac12 b}\\
&.\frac{\sin(A+\omega).\sin(A-\frac12\omega)^2}{\sin A^3}.\frac{\sin B^3}{\sin(B+\omega).\sin(B-\frac12\omega)^2}\Biggr)
\end{aligned}
\]
\noindent oder
\[
\log\frac{a\sin B}{b\sin A}=\frac{b^4-a^4}{180}+.\, .\, .+\frac{\omega\omega}{4}\left(\frac{1}{\sin B^2}-\frac{1}{\sin A^2}\right).
\]
\noindent Auch ist hinreichend scharf
